Hermite gave me bounded evidence about action-representation fidelity in an offline codec benchmark. Separately, CR5 work showed how timestamp alignment, schema checks, and motion validation shape the data entering a learning pipeline. These independent foundations lead to my next question: how can efficient processing and scalable learning preserve meaningful trajectory information as embodied datasets grow? The MSc in Artificial Intelligence and Big Data Computing at The Hong Kong Polytechnic University (PolyU) directly addresses that transition.

I would follow the data path from processing to learning to application. DSAI5204 Efficient Data Processing provides the processing layer for growing trajectory datasets; DSAI5207 Modern Deep Learning strengthens the models that learn from those representations; and DSAI5201 Artificial Intelligence and Big Data Computing in Practice places both within an applied workflow. Their availability remains subject to curriculum finalisation and departmental discretion. If I met the published eligibility conditions and selected the dissertation route, the nine-credit DSAI5902 Dissertation for Artificial Intelligence \& Big Data Computing could support a study asking how processing choices affect action-representation fidelity and downstream behavior as trajectory data scale. I could contribute experience from tokenizer evaluation and real-robot data-quality checks, then carry the resulting methods into scalable embodied-AI R\&D.
